\section{Evaluation}\label{sec:eval}

We ask seven questions. \textbf{RQ1:} Does compiling against the data find mistakes before any data
is processed? \textbf{RQ2:} What does data-aware compilation cost and gain? \textbf{RQ3:} Can a
compiled program be reused on edited batches? \textbf{RQ4:} On real data that changes from month to
month, compared with data-validation tools? \textbf{RQ5:} For real pandas scripts? \textbf{RQ6:} Do
certificates compose across stages? \textbf{RQ7:} What does reuse cost at run time? We also report
how we tested soundness and the bugs the evaluation found in our own implementation.

\paragraph{Setup.} Intel Core i5-12450H (8 cores), 16\,GB RAM, Windows~11; MSVC 14.44 (Release);
Python 3.14, pandas 3.0, NumPy 2.5, OpenCV 5.0, pyarrow 25. pandera 0.34 and Great Expectations
0.18.22, whose onboarding data assistant builds a suite from a sample (we found no such assistant in
the 1.x line), run in a Python 3.11 environment, TFDV 1.21 in WSL (Ubuntu, Python 3.12), since it has
no Windows build. Times include starting Python.

\paragraph{Scoring reuse.} RQ3--RQ6 score every way of running a new batch $T$ against the ground
truth: a fresh compilation of the pipeline for $T$, run. A policy either runs a program on $T$ or
refuses it. Its outcome is \emph{correct} if it wrote exactly the files of the fresh run or stopped
where the fresh run stops (a refusal, or a failing check, when the fresh compilation fails or its run
stops); \emph{needless} if it refused $T$ although the fresh compilation runs, which costs a
recompilation; and \emph{wrong} if it finished with other files than the fresh run, finished where the
fresh run fails, or stopped where it does not. The policies are: the \emph{fingerprint} of the
compiled file; the compiled file's own facts (types, no missing values, and the range of every column
read, a schema as TFDV would record it with ranges); pandera's \texttt{infer\_schema} (dtypes,
nullability, numeric ranges); Great Expectations' onboarding data assistant (an inferred suite of
column statistics); TFDV's \texttt{infer\_schema} (types, presence, string domains); the \emph{full
certificate}; and \emph{selective deoptimization}. The validators are set up on the file the pipeline
was compiled with, once on every column (as they are used) and once on the columns the program reads
(more favourable to them); when a validator admits $T$, the optimized program runs on it as it is.

\subsection{RQ1: Mistakes found before any data is processed}

\paragraph{Injected mistakes.} We injected 24 mistakes into the five example programs, one per
mutant~\cite{jia2011}: misspelled columns, type errors (normalizing text, comparing a number with
text), uses after removal, broken contracts (range, not-null, foreign key, schema, channels, width,
match rate), a broken image, and three well-typed programs with silently wrong output. \PL reports all
24 at compile time and nothing on the unmodified programs, in \RqOneMsMin--\RqOneMsMax\,ms; a
script-style interpreter of the same programs crashes part-way on 20 and finishes 4 with wrong output
(text filled into an integer column, a threshold on a normalized column that removes every row, and
two pipelines that write black images). On 16 million Parquet rows, a misspelled column is reported in
\ScTabNamePl\,s instead of after \ScTabNameScript\,s, and float pixels saved as PNG in \ScImgFloatPl\,s
for 3,000 images, where the script runs \ScImgFloatScript\,s and writes black images.

\paragraph{Real failures.} We chose 17 error messages that Python libraries print for common
data-preparation failures, searched GitHub issues for each and read the first 12 results in order; an
issue qualified if the failure happened while loading, transforming or saving data files and its root
cause was stated, at most three per message: \RbTotal{} issues~\cite{islam2019,humbatova2020}. \PL
reports \RbEarly{} of them (\RbEarlyPct\%) before any data is processed (\RbCompile{} without
annotation, \RbContract{} once the requirement is declared with \texttt{EXPECT}); \RbDesign{} cannot
occur in a \PL program, and \RbNone{} are not covered (input formats such as HEIC or GeoTIFF,
operations such as pivots, dates, per-image label files). \RbPass{} of \RbCases{} reproductions behave
as classified.

\subsection{RQ2: Cost and gain of data-aware compilation}

Analysing every cell of a one-million-row CSV file takes \CsvMillion\,s, \CsvSpeedup$\times$ faster
than pandas' default reader reads it; Parquet compilation stays at \PqSixteen\,s from one to 16
million rows; reading 3,000 image headers is \HdrRatio$\times$ faster than decoding the images. On a
200,000-row, 40-column table with two contracts and a filter, the optimizer proves both contracts and
loads 5 of 40 columns; the optimized program processes the data in \TabProcOpt\,s instead of
\TabProcOz\,s (\TabCut\% less) with byte-identical output. Image pipelines stream one image at a
time, in \ImgPlMem\,MiB of memory instead of \ImgBaseMem\,MiB for loading all 1,000 images.

\subsection{RQ3: Reusing a compiled program on edited batches}\label{sec:rq-edits}

Seven workloads are compiled once: the four table examples, one program run twice on one file, the
Parquet example, images linked to labels, and a sensor pipeline of 20,000 rows per day. Each input
table gets new batches: other days and a warmer season of the sensor process, and, for the examples,
rows resampled with replacement, values moved by noise of 0.1, 0.5 and 1.5 standard deviations, 20\%
missing values in one column, and one value far out of range (\CertBatches{} batches).

\begin{table}[t]
\caption{New batches admitted by each policy (RQ7)}
\label{tab:rq7}
\centering
\footnotesize
\begin{tabular}{@{}lrrrr@{}}
\toprule
Workload & Batches & File facts & Certificate & Ends like fresh \\
\midrule
daily & 8 & 0 & 8 & 8 \\
students & 24 & 7 & 14 & 14 \\
optimize & 24 & 7 & 15 & 15 \\
ml\_split & 48 & 15 & 42 & 42 \\
tableset & 96 & 41 & 96 & 96 \\
parquet & 24 & 10 & 14 & 14 \\
linked\_images & 20 & 16 & 20 & 20 \\
\midrule
all & 244 & 96 & 209 & 209 \\
\bottomrule
\end{tabular}
\end{table}


Both certificates handle every batch correctly (Table~\ref{tab:rq-edits}): the full certificate by
refusing exactly the batches whose fresh compilation fails, selective deoptimization by running them
and stopping at the re-enabled check (\CertDeopted{} batches ran with removals re-enabled). The
fingerprint refuses every edited batch, and so do the compiled file's own facts for most. Gating the
optimized program with the file's own facts is also unsound (\CertFactsUnsound{} wrong runs): on the
table set and on the linked images, resampled and drifted rows repeat keys or lose matches, so a
contract the compiler had proven by reading the values (\texttt{UNIQUE}, \texttt{REFERENCES},
\texttt{MATCH\_RATE}) no longer holds and the fresh compilation reports it, while the file's facts,
which are facts of single columns, admit the batch and the optimized program finishes. \PL runs such
checks again whenever an input changed.

\paragraph{Image folders.} A pipeline that makes images gray and $64\times64$, compiled on a folder
of gray $64\times64$ images, has its \texttt{GRAYSCALE} and \texttt{RESIZE} removed speculatively. We
gave it \ImgBatches{} new folders: gray $64\times64$ images, some color images, other sizes, both, a
file that is not an image, and 16-bit images (five each), and compared the written images byte by
byte with a fresh compilation. Selective deoptimization handles \ImgBDeoptCorrect{} correctly,
re-running \texttt{GRAYSCALE} or \texttt{RESIZE} on \ImgBDeopted{} folders; the full certificate
refuses those folders, and the optimized program run on them without a check writes color or
wrongly sized images (\ImgBTrustedUnsound{} folders). The \ImgBDeoptNeedless{} needless refusals are
the 16-bit folders: bit depths are fixed in every certificate because they decide the pixel type the
generated code assumes, even where, as here, no instruction depends on it; certifying them would need
the code generator's assumptions as obligations too.

\subsection{RQ4: Real monthly data, against data-validation tools}\label{sec:rq-real}

Three public datasets come in monthly batches: hourly traffic and weather on a Minnesota interstate
(UCI, 2012--2018, 63 months of CSV), daily bike rentals in Washington (UCI, 2011--2012, 24 months),
and New York green-taxi trips (NYC TLC, 2022--2024, the 36 Parquet files as published). For each we
wrote one cleaning pipeline as a data engineer would: range contracts on the measured quantities,
\texttt{NOT\_NULL} on keys, fills, a filter that removes impossible readings, a normalization and a
selection. Each pipeline is deployed twice: compiled on the first month and run on every later one,
and compiled on the month before each month. Batches are cut from the raw files, so no value is
reformatted.

\begin{table}[t]
\caption{Real monthly batches (RQ8). Facts and Cert.: months admitted by the compiled month's own facts and by the certificate; Fresh: admitted months that end like a fresh compilation; Refused: refused although a fresh compilation runs. Traffic and Bikes are UCI CSV files, Taxi is NYC TLC green-taxi Parquet.}
\label{tab:rq8}
\centering
\footnotesize
\begin{tabular}{@{}llrrrrr@{}}
\toprule
Data & Anchor & Months & Facts & Cert. & Fresh & Refused \\
\midrule
Traffic & first & 62 & 0 & 59 & 59 & 1 \\
 & previous & 60 & 2 & 58 & 58 & 1 \\
Bikes & first & 23 & 0 & 22 & 22 & 1 \\
 & previous & 23 & 0 & 22 & 22 & 1 \\
Taxi & first & 35 & 0 & 11 & 11 & 17 \\
 & previous & 35 & 0 & 33 & 33 & 2 \\
\bottomrule
\end{tabular}
\end{table}

\begin{figure}[t]
\centering
\begin{tikzpicture}
\begin{axis}[
    xbar stacked, bar width=7pt, width=\columnwidth, height=6.2cm,
    xmin=0, xmax=\RealRuns, enlarge y limits=0.06,
    xlabel={real monthly batches}, xlabel style={font=\scriptsize},
    symbolic y coords={TFDV read,TFDV all,GX read,GX all,pandera read,pandera all,File facts,Fingerprint,Full cert.,Deopt.},
    ytick=data, yticklabel style={font=\scriptsize}, xticklabel style={font=\scriptsize},
    legend style={font=\scriptsize, at={(0.5,1.02)}, anchor=south, legend columns=3, draw=none},
    reverse legend=false]
\addplot[fill=black!70, draw=black!70] coordinates {
    (\RealTfdvReadCorrect,TFDV read) (\RealTfdvAllCorrect,TFDV all) (\RealGxReadCorrect,GX read)
    (\RealGxAllCorrect,GX all) (\RealPanderaReadCorrect,pandera read) (\RealPanderaAllCorrect,pandera all)
    (\RealFactsCorrect,File facts) (\RealFpCorrect,Fingerprint) (\RealFullCorrect,Full cert.)
    (\RealDeoptCorrect,Deopt.)};
\addplot[fill=black!20, draw=black!40] coordinates {
    (\RealTfdvReadNeedless,TFDV read) (\RealTfdvAllNeedless,TFDV all) (\RealGxReadNeedless,GX read)
    (\RealGxAllNeedless,GX all) (\RealPanderaReadNeedless,pandera read) (\RealPanderaAllNeedless,pandera all)
    (\RealFactsNeedless,File facts) (\RealFpNeedless,Fingerprint) (\RealFullNeedless,Full cert.)
    (\RealDeoptNeedless,Deopt.)};
\addplot[pattern=north east lines, pattern color=black, draw=black] coordinates {
    (\RealTfdvReadUnsound,TFDV read) (\RealTfdvAllUnsound,TFDV all) (\RealGxReadUnsound,GX read)
    (\RealGxAllUnsound,GX all) (\RealPanderaReadUnsound,pandera read) (\RealPanderaAllUnsound,pandera all)
    (\RealFactsUnsound,File facts) (\RealFpUnsound,Fingerprint) (\RealFullUnsound,Full cert.)
    (\RealDeoptUnsound,Deopt.)};
\legend{correct, refused needlessly, wrong}
\end{axis}
\end{tikzpicture}
\caption{RQ4: \RealRuns{} real monthly batches, by admission policy. \emph{Correct}: ran and wrote
the fresh compilation's files, or stopped where it stops; \emph{refused needlessly}: a recompilation
was needed although the fresh compilation runs; \emph{wrong}: other files, or a run where the fresh
compilation fails. Validators gate the optimized program; \emph{all}/\emph{read}: set up on every
column, or on the columns the program reads.}
\label{fig:policies}
\end{figure}


Over \RealRuns{} runs (Table~\ref{tab:rq-real}, Fig.~\ref{fig:policies}), selective deoptimization
handles \RealDeoptCorrect{} (\RealDeoptCorrectPct\%) exactly as a fresh compilation would, with
\RealDeoptUnsound{} wrong; the full certificate handles \RealFullCorrect{}. The difference is
\RealRescued{} months the full certificate refuses. In three, the traffic sensor reports temperatures
of 0\,K; the re-enabled range check stops the run, as the fresh compilation fails. In four, a month has
an hour with 9,831\,mm of rain or a day with humidity 0, which the compiled month had not; the removed
cleaning filter runs again and the run writes what the fresh compilation writes. In two, a taxi trip
has a fare of \$2,020; the re-enabled check runs after the filter that removes that trip, passes, and
the run finishes. All \RealDeoptNeedless{} remaining refusals are one event: in February 2023 the
taxi publisher changed \texttt{passenger\_count} and \texttt{payment\_type} from double to integer,
and the program compiled on January 2022 reads them as \textsc{Float}. A type is a fact no run-time
code of the compiled program stands in for; monthly compilation follows the change after one month.

The validators are either strict or unsound. Great Expectations' profile constrains means, medians,
quantiles and unique counts, and pandera's inferred schema the compiled month's range of every
column; set up on every column or only on the columns read, they refuse
\RealValidatorStopMin--\RealValidatorStopMax\% of the months although a fresh compilation runs them.
TFDV infers types, presence and string domains but no numeric ranges. On every column it refuses
\RealTfdvAllNeedlessPct\% of the months, for string values of columns the program never reads (a new
holiday name refuses a traffic month); on the columns read it refuses fewer, but both setups admit
months on which the optimized program writes other files than a fresh compilation
(\RealTfdvReadUnsound{} months): the taxi months after the type change, written with the old column
types, and traffic months with an impossible rain reading that a removed filter should have dropped.
Generic validators do not know which facts the compiled program relies on.

\subsection{RQ5: Real pandas scripts}\label{sec:rq-corpus}

\paragraph{Corpora.} GitHub code search for Python files containing \texttt{read\_csv},
\texttt{fillna} and \texttt{Embarked} (the Kaggle Titanic data) returned \CorpusAll{} distinct files,
and for \texttt{read\_csv}, \texttt{fillna} and \texttt{LotFrontage} (Kaggle House Prices) \HousesAll{}
files. We used the first 100 Titanic scripts as a development set, to choose which idioms the lifter
handles; the remaining \CorpusTest{} and all House Prices scripts are held out: we did not look at them
before the final run, and the lifter was not changed for House Prices. A script is in scope if it reads
the training file with \texttt{pd.read\_csv(path)} at module level. The House Prices files mark missing
values with the text \texttt{NA}, which pandas reads as missing; we give both sides the same files with
empty cells instead, \PL's notation.

\paragraph{Protocol.} The lifted program is compiled against the training file and given eleven new
batches: Kaggle's real test file (no target column, and missing values in columns where the training
file has none), five bootstrap resamples of the training rows, and five copies with 20\% extra missing
values in one column. Each batch is run by the compiled program and, as ground truth, by the lifted
pandas statements; the cleaned tables are compared value by value.

\begin{table}[t]
\caption{Real pandas scripts from GitHub (RQ5): held-out scripts, lifted, compiled once on the training file and run on new batches; the ground truth is the lifted pandas statements run on the same batch.}
\label{tab:rq-corpus}
\centering
\footnotesize
\begin{tabular}{@{}lr@{}}
\toprule
& Titanic \\
\midrule
Held-out scripts & 272 \\
Lifted, in scope & 130 \\
Compiled & 124 \\
Runs on new batches & 1364 \\
\emph{Full certificate:} admitted & 1347 \\
\quad identical to the script & 1347 \\
\quad refused, script runs & 10 \\
\emph{With deoptimization:} admitted & 1357 \\
\quad identical to the script & 1357 \\
\quad stopped by a check & 0 \\
\quad different (unsound) & 0 \\
\quad refused, script runs & 0 \\
\quad ran with removals re-enabled & 10 \\
\bottomrule
\end{tabular}
\end{table}


\paragraph{Results.} With deoptimization, \CorpusSame{} Titanic runs and \HousesSame{} House Prices
runs reproduce the scripts' results exactly, and none differ (Table~\ref{tab:rq-corpus}). The
difference to the full certificate is the price of optimizing for the compiled data: the training
files have no missing \texttt{Fare} (Titanic) or \texttt{MSZoning}, \texttt{KitchenQual} and other
columns (House Prices), so the compiler removes the scripts' \texttt{fillna} on them, and the full
certificate then refuses every batch with missing values there, including Kaggle's own test files.
Deoptimization runs the removed \texttt{fillna} again for exactly those batches.

\subsection{RQ6: Composing stages}\label{sec:rq-compose}

Each workload of RQ4 gets a second stage that reads the cleaned table, as a feature step would: range
and not-null contracts, a fill, a filter and a normalization. Both stages are compiled once (stage~2 on
what stage~1 writes for the anchor month), and \texttt{tools/contracts.py} compares stage~1's guarantee
with what stage~2 relies on. Every later month runs through both compiled stages, and stage~2 is also
compiled and run fresh on stage~1's output.

\begin{table}[t]\caption{Results pending.}\label{tab:rq-compose}\end{table}


Across \ComposeMonths{} months through both stages, no fact the composition reported as guaranteed ever
made stage~2 refuse or deoptimize (Table~\ref{tab:rq-compose}). The facts it reported as open are
real: stage~1 never constrains, for example, taxi tips, which stage~2 expects in $[0,500]$. With the
\texttt{EXPECT} statements the tool suggested added to stage~1, every fact composes, and a month that
breaks one now stops in stage~1, where the data enters, instead of deoptimizing downstream.

\subsection{RQ7: The cost of reuse}\label{sec:rq-cost}

Reusing a compiled program costs its admission check and, with deoptimization, the facts of the columns
item certificates mention. Over the \OhMonths{} real months on which every variant finished, the
median run took \OhTrusted\,s with no check at all, \OhFull\,s with admission against the full
certificate, and \OhDeopt\,s with admission and deoptimization (\OhPct\% over no check); compiling the
month afresh took \OhCompile\,s before running, and compiling and running \OhFresh\,s. Deoptimization
needs one pass over a few columns that the program reads anyway; recompiling reads every cell again.

\subsection{Testing soundness}\label{sec:soundtest}

The theorems rest on monotonicity and soundness of the analysis, which we test rather than prove for
the implementation. Randomized property tests generate tables (with missing, out-of-range, pruned and
dropped columns) and image folders, decide admission and deoptimization with the certificates, and
compile each input afresh: every input the full certificate admits keeps every proof, and every input
admitted with deoptimization compiles without new errors and keeps every removal it does not
deoptimize (or stops at a re-enabled check exactly where the fresh compilation reports a violation).
For the running example, \RandFull{} of 400 random tables fit the full certificate and \RandDeopt{} are
admitted with deoptimization; \RandFolders{} random image folders are admitted with deoptimization,
all keeping their removals.

\paragraph{What the evaluation found in our implementation.} Each of the following was a case where
the implementation did not follow the definitions; each is fixed and has a regression test. (1)~The
run-time \texttt{MATCH\_RATE} check, always eliminated before certificates existed, ran for the first
time on changed labels and counted broken images that \texttt{REMOVE\_CORRUPT} had removed. (2)~The
first certificate covered only the columns the optimized program loads, so a proven \texttt{NOT\_NULL}
on a column pruned afterwards was unprotected. (3)~Admission let \textsc{Integer} data into a
\textsc{Float} column and an empty column into any type; the property test found inputs whose fresh
compilation decided differently. (4)~On the held-out Titanic set, scripts drop \texttt{Survived} before
any use; with that \texttt{DROP} removed, admission treated the column as optional and admitted
Kaggle's test file, on which the script fails. (5)~The lifter translated a \texttt{drop} that pandas
rejects.

\subsection{Threats to validity}

The error messages of RQ1 were chosen by the author and one person classified the issues; the
mutants and the script-style interpreter are ours. Performance numbers come from one machine. The
pipelines of RQ4 and RQ6 are ours, written once before running them; the months are not. Validators are
used with their automatic profilers, as a new user would; a hand-tuned suite can do better, but then
encodes what the certificate derives. RQ5 compares with the lifted statements, not the whole scripts,
and the lifted prefixes are short. Metadata-only soundness assumes truthful Parquet statistics. The
mechanized theorem assumes monotonicity and soundness of the analysis, which we test. The language has
no joins, loops or user functions, and certificates are facts about input columns: a fact needed only
after a filter is still required of every input row.
